\section{Open-weight coding stream and same-task shadow observations}
\label{app:open_coding}
This laboratory study collected new open-weight model outputs under a
frozen exposure schedule. It is neither a production-user A/B test nor a
reassignment of an existing complete outcome table. The original online
monitor ran during collection, without stopping execution. The
normal-mixture reanalysis below was chosen after the outcomes were known;
we distinguish physical collection from post-hoc inferential choices.

\paragraph{Roster, model and workflows.}
The complete roster contains 427 sanitized MBPP tasks
\citep{austin2021program} and 164 HumanEval tasks \citep{chen2021evaluating}.
We used the MLX four-bit conversion of Qwen2.5-Coder-7B-Instruct,
revision \texttt{019cc73c45c770444708a6dd8690c66243cc5c80}, with
\texttt{mlx\_lm} 0.31.3, Python 3.12.13 and one Apple M5 machine with
10 CPU cores and 32 GiB memory. The snapshot inventory contains
4,295,890,004 bytes, including tokenizer and metadata files; this is not
a memory-usage measurement. Sampling used temperature 0.7, top-$p$ 0.95
and a 1,024-token completion cap per call. A generated one solution. B generated
a solution and its own tests, executed those tests, and made up to two
repairs from execution feedback, using at most four model calls.
Hidden verification tests were not supplied for repair. MBPP prompts used
the task description and reference-derived function signature; standard
HumanEval prompts include their original illustrative doctests.
The exact collection definitions, sampling configuration and pinned
benchmark/model identities accompany the analytical data.

\paragraph{Assignment and observation units.}
Seed 20260918 generated a permutation of all 591 tasks and an AB/BA
orientation for each consecutive pair. The first pass exposed each
arrival once to its assigned workflow, one execution at a time. The
second pass ran the complementary workflow on every task in the same
order. There are 295 complete first-pass pairs and one prespecified
unpaired arrival; the latter contributes only to the 591 same-task
comparisons. Both passes finished, giving 1,182 episodes without missing
enrolled tasks. The optional second trial was disabled. This is one
collected schedule, not independent replications of the task roster.
The design's pseudorandom seed and full assignment array are retained;
neither random order nor distinct task identifiers proves independent
episode outcomes or stationary execution latency.

\paragraph{Endpoints and comparison.}
Success requires the hidden checks to exit successfully and the verifier's
sentinel to be observed. Timeouts and failed checks are unsuccessful.
The hierarchy compares success first, measured workflow latency second,
and completion-token count third. For each resource tier, a difference
is decisive only if it exceeds 10\% of the larger recorded value.
Resource tiers apply only after joint success; joint failure ties.
Positive scores favor B. Workflow latency includes model calls and
agent self-tests and excludes hidden verification and model loading.
It is local elapsed workflow time, not production queueing or calendar
time to a monitoring decision. Token counts include all calls within
each workflow. No token counts were estimated; there were no recorded
connection errors or retries. Two A episodes and no B episodes had
verifier timeouts. Execution limits and heuristic flags do not establish
adversarial robustness or correctness beyond these archived tests.

\begin{table}[h]
\centering\small
\caption{Observed resources across all 591 tasks per workflow. A is
single shot; B uses self-tests and repair. Hidden verification is excluded
from workflow latency and the self-test execution count.}
\label{tab:open_coding_resources}
\begin{tabular}{lrr}\toprule
 & A & B\\\midrule
Successful tasks & 433 & 433 \\
Model calls & 591 & 1,748 \\
Prompt tokens & 71,218 & 528,191 \\
Completion tokens & 41,491 & 185,088 \\
Total tokens & 112,709 & 713,279 \\
Mean workflow latency (s) & 2.800 & 12.476 \\
Median workflow latency (s) & 2.079 & 10.131 \\
Self-test executions & 0 & 1,157 \\
\bottomrule\end{tabular}

\end{table}

\paragraph{Observed first-pass contrast and post-hoc running-mean analysis.}
Among the 295 first-pass pairs, B has 69 wins, 18 ties and 208 losses,
giving net benefit $-139/295=-0.4712$. Success decides 67 wins and 52
losses, latency decides two wins and 156 losses, and the token tier
decides no pair. Thus the success difference is $15/295=0.05085$,
whereas composite preference favors A under the declared hierarchy.
The original hypotheses of higher B success and positive net benefit
were not supported by the complete same-task success counts and this
first-pass preference result, respectively.

Let $Z_k$ be the observed first-pass sign and $D_k$ its success difference.
For a filtration $\mathcal F_k$ containing the fixed roster/schedule and
revealed records through pair $k$, but no future outcomes, define
$\mu_k=\mathbb E[Z_k\mid\mathcal F_{k-1}]$ and
$\nu_k=\mathbb E[D_k\mid\mathcal F_{k-1}]$. We target the running
averages $n^{-1}\sum_{k\leq n}\mu_k$ and
$n^{-1}\sum_{k\leq n}\nu_k$. These history-conditional means allow
task composition, machine state and serial dependence. Conditioning on
the full assignment schedule does not simultaneously give a conditional
fair-coin causal interpretation. Such an interpretation instead requires
orientations to be unrevealed and nonanticipating at their pairs, and
appropriate history-dependent assignment-response laws.

For each process separately, Theorem~\ref{thm:normal_cs} gives the fixed
normal-mixture construction with $V_n=n$, $\rho=100$ and $\alpha=0.05$:
\[
r_n=\frac{\sqrt{(n+100)\log\{(n+100)/(100\cdot0.05^2)\}}}{n}.
\]
Apply $[\overline Z_n-r_n,\overline Z_n+r_n]\cap[-1,1]$ and its
success-difference analogue, without intersections across times because
the targets can move. At $n=295$, $r_n=0.18284$, yielding the marginal
bands $[-0.65403,-0.28835]$ for net benefit and
$[-0.13199,0.23369]$ for success difference. The first negative
net-benefit upper endpoint occurs at pair 60. The success lower endpoint
never exceeds $-0.03$: neither success noninferiority nor a guarded
deployment is certified. These two separate 95\% constructions do not
form a joint 95\% region. The theorem covers a specified fixed procedure;
the post-hoc choice of this analysis is not adjusted for selection.
The crossing is a chronological reanalysis, produced after all planned
episodes completed, and implies no realized savings.

\paragraph{Same-task descriptions and generalization limits.}
Both workflows succeed on 433/591 tasks: 393 joint successes, 40
successes unique to each workflow, and 118 joint failures. The
same-task hierarchy gives 41 B wins, 121 ties and 429 losses, hence
$-388/591=-0.6565$ net benefit. This is descriptive: shared pair
orientations and complementary passes can create dependence and period
effects. We attach no contributed task-level interval, sampling error for
the difference between contrasts, or task-superpopulation claim.
The running-mean bands above do not cover the all-pairs fixed-roster
functional. One model, one machine and one schedule cannot establish
hardware-invariant workflow ranking or transfer to production tasks.
Observed B/A ratios are 4.46 for mean workflow latency, 4.46 for
completion tokens and 6.33 for total tokens; these are distinct endpoints,
not a monetary-cost or energy estimate.

\paragraph{Reproduction and retained provenance.}
The supplement includes a row-preserving projection of every episode's
metrics, all task assignments, original source hashes, model/tokenizer
file hashes and pinned benchmark URLs. It omits generated programs,
self-tests, verifier stderr and identifying local paths. Reproduction
recalculates every comparison and marginal band from the archived
verifier labels and resources; it does not regenerate model outputs or
rerun the hidden tests. Collection definitions are retained as
nonexecuted historical text. Superseded contributed E2/R2 uncertainty and
betting-capital endpoint routines are excluded from this analysis.
